%% tags: [breadth first search]
%% source: 2023-fa-midterm_02
\begin{prob}

    Consider the code below, which is a modification of the BFS algorithm
    presented in lecture: there are two new lines, marked with comments.

    \begin{minted}[autogobble]{python}

        from collections import deque

        def modified_bfs(graph, source):
            """Start a BFS at `source`."""
            status = {node: 'undiscovered' for node in graph.nodes}
            status[source] = 'pending'
            pending = deque([source])

            while pending:
                u = pending.popleft()
                for v in graph.neighbors(u):

                    for node in graph.nodes:  # < -- this is the modification
                        print(status[node])   # < -- this is the modification

                    if status[v] == 'undiscovered':
                        status[v] = 'pending'
                        pending.append(v)

                status[u] = 'visited'
    \end{minted}

    What is the time complexity of this algorithm, in terms of $|V|$ and $|E|$?
    You may assume the graph is undirected and connected, so that all nodes are
    reachable from the source.

    \begin{choices}
        \choice $\Theta(1)$
        \choice $\Theta(V)$
        \choice $\Theta(E)$
        \choice $\Theta(V + E)$
        \correctchoice $\Theta(V E + V)$
    \end{choices}

    \begin{soln}
        Because the graph is connected, $|E| \geq |V| - 1$, and so $\Theta(VE)$ and $\Theta(VE + V)$ are equivalent. Only $\Theta(VE + V)$ is an option so $\Theta(VE + V)$ is the correct answer.
    \end{soln}

\end{prob}
